\section{Introducción: por qué este episodio es una nota de ruta para emprender con aplicaciones de IA}

Esta sección sitúa primero el episodio. EP95 es una «entrevista en relevo»: la primera parte se grabó en el otoño de 2024, cuando Xiao Hong (肖弘) acababa de cerrar una ronda de financiamiento, y trata sobre el emprendimiento en serie, el capital, Monica y la metodología de las aplicaciones de IA; la segunda parte se grabó después del Año Nuevo chino de 2025, cuando DeepSeek había transformado el ecosistema de base de las aplicaciones de IA en China, y Xiao Hong empieza a hablar de Manus, que estaba por lanzarse y que fue el «primer disparo» de los productos Agent en China. Por ello, este episodio no es una historia estática de un fundador, sino un registro continuo de las reflexiones de un emprendedor de aplicaciones de IA en medio de un cambio acelerado en las capacidades de los modelos.

\begin{importantbox}{Tesis central del episodio}
La línea principal de Xiao Hong puede resumirse en tres frases: primero, una empresa de aplicaciones de IA debe anticipar las capacidades de los modelos y los movimientos de las grandes empresas, y preparar con antelación el contenedor de producto; segundo, las capacidades de los modelos tenderán a la igualación tecnológica, pero la barrera de entrada de una empresa de aplicaciones proviene del escenario de uso, la marca, la comprensión de los usuarios, la retroalimentación y la ejecución; tercero, el mundo no es una extrapolación lineal: el Founder debe pensar en términos de teoría de juegos y convertirse en una variable relevante de la situación.
\end{importantbox}

\begin{knowledgebox}{Nota sobre la estrategia visual}
Este video es una entrevista con plano fijo, sin proyección de pantalla ni demostración de producto. El cuerpo de la nota usa la portada solo para identificar la fuente; todas las imágenes del cuerpo son diagramas conceptuales propios que explican la estructura de la entrevista en relevo, el costo del VC, la anticipación de los movimientos de las grandes empresas, Monica, la clasificación de las aplicaciones de IA, el producto Agent de Manus y el pensamiento estratégico del Founder.
\end{knowledgebox}

\begin{figure}[H]
\centering
\includegraphics[width=0.96\textwidth]{xiaohong-relay-interview.png}
\caption{Mapa de la entrevista en relevo con Xiao Hong: las dos entrevistas registran las reflexiones continuas de un emprendedor de aplicaciones de IA en un entorno inestable. Diagrama conceptual propio, elaborado a partir de la estructura de capítulos del video completo.}
\end{figure}

\begin{knowledgebox}{Lectura de la figura: no es una línea recta, sino dos cortes de estado}
La figura va desde después del financiamiento del otoño de 2024, pasando por el primer emprendimiento, Monica y el impacto de DeepSeek en el Año Nuevo chino de 2025, hasta Manus. Al leerla, conviene notar que un mismo Founder enfrenta, en distintos momentos, capacidades externas de los modelos, entornos de capital, oportunidades de producto y problemas de organización diferentes; por eso sus juicios también cambian.
\end{knowledgebox}

\subsection{Ruta de lectura}

Esta sección presenta la ruta de lectura. El episodio completo puede dividirse en seis preguntas: primera, ¿qué lecciones de producto, de monetización y de capital le dejó a Xiao Hong su primer emprendimiento? Segunda, ¿por qué «el VC es un medio muy caro de obtener financiamiento»? Tercera, ¿por qué Monica eligió la extensión de navegador, el mercado internacional y el all-in-one AI assistant? Cuarta, ¿cómo se clasifican las aplicaciones de IA: complemento del escenario principal, cambio en las capacidades de los modelos y desbordamiento desde dominios verticales? Quinta, ¿por qué se entiende a Manus como el primer disparo de los Agent en China? Sexta, ¿por qué el Founder debe pasar del razonamiento lógico al pensamiento de teoría de juegos?

\noindent\begin{tabularx}{\textwidth}{p{0.20\textwidth}p{0.34\textwidth}X}
\toprule
Pregunta de lectura & Material de la entrevista & Juicio que se debe formar \\
\midrule
Qué aprendió del primer emprendimiento & Herramientas para cuentas oficiales de WeChat, SCRM, monetización, venta de la empresa & Un producto de tipo herramienta debe monetizarse pronto, sin dejarse desviar por grandes narrativas. \\
Cómo influye el capital en las decisiones & El VC es caro, el peso de los inversionistas, la dimensión del capital & El financiamiento es una herramienta; no es el jefe ni es la estrategia en sí. \\
Por qué se creó Monica & Extensión, mercado internacional, ChatGPT for Google, retroalimentación de usuarios & Una aplicación de IA debe ocupar el contenedor de producto que surge cuando las capacidades del modelo se desbordan. \\
Cómo encuentra oportunidades una empresa de aplicaciones & Jasper, ChatGPT, Cursor, DeepSeek, Manus & Observar la brecha entre las capacidades del modelo y el escenario principal del usuario. \\
El significado de Manus como producto & Tareas, herramientas, entorno, memoria, entrega de resultados & El Agent pasa de la conversación a un sistema capaz de hacer cosas. \\
Cómo piensa el Founder & No linealidad, teoría de juegos, la edad frente a la época, codicia, ira e ignorancia & Emprender no es un razonamiento lineal; hay que cambiar la situación. \\
\bottomrule
\end{tabularx}

\begin{warningbox}{Límites temporales}
La primera entrevista se realizó en el otoño de 2024. En esa parte, «este año» se refiere casi siempre a 2024 y «el año pasado», a 2023. La segunda entrevista se realizó después del Año Nuevo chino de 2025 y trata la situación posterior al auge de DeepSeek y previa al lanzamiento de Manus.
\end{warningbox}

\subsection{Resumen de la sección}

EP95 es una nota de metodología para emprender con aplicaciones de IA. Puede leerse junto con la última entrevista de EP128 sobre Peak/Manus, EP101 sobre YouWare y EP110 sobre el informe técnico de Agent: aquella describe la organización de Manus en una etapa posterior y su situación antes de la adquisición, mientras que este episodio recoge las reflexiones del fundador antes del lanzamiento de Manus y la metodología de base para emprender con aplicaciones.
